\begin{frame}[fragile]{Web 1.0 - Static HTML}

\begin{itemize}
	\item HTML was invented to display Document Object Model (DOM) information on web browsers. 
	\item When web servers get a request from clients (mostly web browsers), they send HTML files as a response (Web 1.0). 
	\item CSS (Cascading Style Sheet) was invented to render HTML elements on a web browser more flexibly. 
\end{itemize}

\begin{columns}%
\begin{column}{0.5\textwidth}%
\begin{center}

\includegraphics[width=0.3\textwidth]{topics/SE_programming/javascript/pic/history/html.png}
\end{center}
\end{column}%
\begin{column}{0.5\textwidth}%
\begin{center}

\includegraphics[width=0.22\textwidth]{topics/SE_programming/javascript/pic/history/css.png} 
\end{center}
\end{column}%
\end{columns}%
\end{frame}

\begin{frame}[fragile]{Web 2.0 - CGI}

\begin{itemize}
	\item With the invention of CGI (Common Gateway Interface) (Web 2.0), \underbar{servers} can generate HTML DOM dynamically from the request. 
	\item JavaScript was invented to manipulate HTML structure (DOM) on the \underbar{clients'} side. 
\end{itemize}

\begin{columns}%
\begin{column}{0.5\textwidth}%
\begin{center}

\includegraphics[width=0.3\textwidth]{topics/SE_programming/javascript/pic/history/js.png}
\end{center} 
\end{column}%
\begin{column}{0.5\textwidth}%
\begin{center}
\includegraphics[width=0.7\textwidth]{topics/SE_programming/javascript/pic/history/dom.png}
\end{center} 
\end{column}%
\end{columns}%

\end{frame}

\begin{frame}[fragile]{}

\begin{itemize}
	\item With the invention of JavaScript, we can have a dynamic web page (DOM) both from a server-side (CGI) and a client-side (JavaScript). 
	\item This code shows JavaScript code that accesses an element with ``one'' id, for example, \verb|<h1 id="one">...</h1>|, and displays the content using the alert() method. 
\end{itemize}

\begin{Code}{}%
\begin{lstlisting}[firstnumber=1, label=glabels, xleftmargin=0pt, basicstyle=\scriptsize]% 

<h1 id="one">Hello</h1> <!-- HTML --> 

<script type="text/javascript"> <!-- JavaScript -->
  var text = document.getElementById("one").innerHTML;
  alert(text); <!-- displays Hello -->
</script>
\end{lstlisting}%   
\end{Code}%

\end{frame}

\begin{frame}[fragile]{}

\begin{itemize}
	\item JavaScript was not invented as an OOP language but as a Functional Programming/Prototype language. 
	\item It is because JavaScript was invented in 10 days by Brendan Eich for the Netscape web browser. 
	\item He had to make the language as simple as possible, not using the complex OOP model. 
\end{itemize}

\begin{columns}%
\begin{column}{0.5\textwidth}%
\begin{center}

\includegraphics[width=0.5\textwidth]{topics/SE_programming/javascript/pic/history/brendan.jpeg}
\end{center} 
\end{column}% 
\begin{column}{0.5\textwidth}%
\begin{center}

\includegraphics[width=0.7\textwidth]{topics/SE_programming/javascript/pic/history/netscape.png}
\end{center} 
\end{column}% 
\end{columns}%

\end{frame}

\begin{frame}[fragile]{}

\begin{itemize}
	\item JavaScript models LISP/Scheme (the first functional programming language). So, JavaScript uses functional programming approaches a lot. 
	\item JavaScript also has the Prototype language feature inherited from Self. 
	\item It has no connection with Java, but to enhance sales of JavaScript, they intentionally named it JavaScript. 
\end{itemize}

\begin{columns}%
\begin{column}{0.5\textwidth}%
\begin{center}

\includegraphics[width=0.5\textwidth]{topics/SE_programming/javascript/pic/history/scheme.jpeg}
\end{center} 
\end{column}% 
\begin{column}{0.5\textwidth}%
\begin{center}

\includegraphics[width=0.6\textwidth]{topics/SE_programming/javascript/pic/history/self.png}
\end{center} 
\end{column}% 
\end{columns}%

\end{frame}
